% Vista científica exacta materializada desde una rama condicional.
% Fuente: source/demostraciones_primarias/14_05b_extension_superviviente_caracter.tex; rama: HMTIncludeCharacterBlock.
% No modifica el contenido matemático de la rama de origen.
\chapter{Carácter positivo de la firma y composición de registros}
\section{Carácter formal universal}

Sea \(\mathcal P=\RR^5\), con coordenadas
\[
 X=(X_A,X_C,X_4,X_{120},X_{270}).
\]
Para \(v=(n_A,s_C,k,\nu_{120},\nu_{270})\in\ZZ^5\), definimos
\[
 \ell_{\rm form}(v;X)
 =n_AX_A+s_CX_C+kX_4
  +\nu_{120}X_{120}+\nu_{270}X_{270},
\]
\[
 \chi_{\rm form}(v;X)=\exp\!\bigl(\ell_{\rm form}(v;X)\bigr).
\]
Así, \(\chi_{\rm form}\) toma valores en el grupo multiplicativo
\(\operatorname{Fun}(\mathcal P,\RR_{>0})\), con producto puntual. Todavía no
ha seleccionado coeficientes globales ni una especie.

\begin{theorem}[Carácter universal de la firma]
\label{thm:caracter-formal-ruta}
El mapa
\[
 \chi_{\rm form}:(\ZZ^5,+)\longrightarrow
 \bigl(\operatorname{Fun}(\mathcal P,\RR_{>0}),\cdot\bigr)
\]
es un homomorfismo. La composición
\[
 \chi_{\rm route}^{\rm form}
 :=\chi_{\rm form}\circ L_{\rm tick}
\]
evalúa formalmente cualquier registro enriquecido y no introduce una unidad
física.
\end{theorem}

\begin{proof}
Para todo \(X\in\mathcal P\), el exponente
\(\ell_{\rm form}(\,\cdot\,;X)\) es aditivo sobre \(\ZZ^5\). Por tanto,
\[
\chi_{\rm form}(v+w;X)
=\chi_{\rm form}(v;X)\chi_{\rm form}(w;X),
\]
y la identidad vale como igualdad de funciones positivas.
\end{proof}

El carácter formal precede a la construcción de los coeficientes globales.
En esta sección no se define todavía un carácter numérico de masa ni una
carta de unidades.

\section{Sección interna mínima sobre las 81 celdas}

El atlas finito contiene 81 celdas orientadas. Al proyectarlas sobre la carta
de familia de ruta se obtienen 56 valores. Denotemos esa proyección por
\[
 \mathcal R:\mathcal C_9\longrightarrow\mathcal R_{56}.
\]
El histograma exacto de sus fibras es
\[
 41\times1,
 \qquad10\times2,
 \qquad2\times3,
 \qquad1\times4,
 \qquad2\times5,
\]
y su suma es 81.

Fijados el origen y el orden lexicográfico de la carta, sea
\(\kappa(x)\in\{0,1,2,3,4\}\) el rango de \(x\) dentro de su fibra.

\begin{theorem}[Selector interno mínimo]
\label{thm:selector-interno-minimo}
La aplicación
\[
 x\longmapsto\bigl(\mathcal R(x),\kappa(x)\bigr)
\]
es inyectiva sobre las 81 celdas. Ningún alfabeto de memoria con menos de
cinco símbolos puede levantar \(\mathcal R\) inyectivamente sobre todo el
atlas.
\end{theorem}

\begin{proof}
Dentro de cada fibra, \(\kappa\) asigna rangos distintos; entre fibras, la
primera coordenada ya es distinta. Existe una fibra de tamaño cinco, así que
el principio del palomar obliga a toda elevación inyectiva a disponer de al menos
cinco valores de memoria. El rango \(0,\ldots,4\) alcanza esa cota.
\end{proof}

La canonicidad es relativa al origen y al orden declarados. \(\kappa\) no es
una carga, un sabor ni una masa. Las trece clases gruesas del atlas producen
multisecciones finitas; sólo la clase central posee el punto fijo \((5,5)\)
bajo la inversión central. Las otras doce requieren información orientada
para seleccionar un punto y no admiten un selector puntual equivariante desde
la clase gruesa sola.

\subsection{Inversión de Möbius y caracteres de ida y retorno}

Sobre la palabra de signos dodecafásica definimos
\[
 C_M(\tau)_m=-\tau_{11-m}.
\]
Esta operación combina la inversión del orden de lectura con el cambio de
hoja. Un carácter lineal
\(\ell_a(\tau)=\sum_ma_m\tau_m\), con
\(a_m=a_{11-m}\), satisface
\[
 \ell_a(C_M\tau)=-\ell_a(\tau),
\]
mientras que toda forma cuadrática simétrica compatible con la reversión
satisface
\[
 Q(C_M\tau)=Q(\tau).
\]
Así, las coordenadas lineales distinguen orientación de ida y retorno y las
cuadráticas conservan el contenido de correlación. Esta paridad explica por
qué \(s_C\) cambia de signo mientras \(\nu_{270}\) puede permanecer
invariante. En el atlas, las siete formas de palabra tienen multiplicidades
\[
 54,2,2,2,7,7,7;
\]
por ello la antihoja se realiza como correspondencia entre fibras y no como
una involución biyectiva obtenida ignorando sus multiplicidades.

\section{Composición orientada de registros antes de la proyección}

Sean \(\Gamma\) y \(\Lambda\) dos historias enriquecidas, \(B\) una
subhistoria común y \(g\) el transporte que identifica fase, hoja,
orientación y posición en la interfaz. En el nivel de registro se define
\[
 \boxed{
 \Gamma\star_B\Lambda
 =\bigl(n_\Gamma+n_\Lambda-n_B,
        e_\Gamma+g e_\Lambda-e_B\bigr).}
\]
La definición presupone también la identificación del borde, los
predecesores torsionales, la corona y la memoria antihoja.

\begin{proposition}[Obstrucción al descenso pentadimensional de la composición de registros]
\label{prop:no-descenso-empalme-cinco}
En general no existe una operación binaria sobre \(\ZZ^5\) que recupere
\(L_{\rm tick}(\Gamma\star_B\Lambda)\) a partir de
\(L_{\rm tick}(\Gamma)\) y \(L_{\rm tick}(\Lambda)\) solamente.
\end{proposition}

\begin{proof}
La operación compone los doce eventos antes de aplicar
\(\operatorname{sgn}\) y las autocorrelaciones. Dos historias pueden tener
la misma cinco-tupla y valores distintos de \(M_5\) o \(A_6\). Al componerlas
con una tercera historia, alguna suma
\(e_a+e_{a+4}+e_{a+8}\) puede cruzar cero en una y no en la otra, cambiando
\(\nu_{120}^{\rm ev}\). Los rombos finitos proporcionan testigos explícitos.
Por tanto, la proyección a cinco enteros no
es una congruencia para \(\star_B\).
\end{proof}

La proposición formaliza una regla operativa esencial: las extensiones se
componen; las firmas se leen después. La partícula compuesta tampoco se
obtiene sumando a ciegas dos filas de una tabla, sino componiendo sus memorias
compatibles y evaluando el resultado.

\section{Dependencia dodecafásica del carácter de persistencia}

Los motores reducidos de período 54 imponen
\[
 e_{m+6}=e_m.
\]
De aquí se siguen las congruencias
\[
 n_A\equiv0\pmod2,
 \qquad
 s_C\equiv0\pmod2,
 \qquad
 \nu_{270}\equiv0\pmod4.
\]
Las firmas de la reconstrucción escalar abreviada violan al menos una de
estas congruencias. Por tanto, su genealogía exige el registro enriquecido y
no procede de ese motor reducido con las reglas declaradas.

Este resultado no es una obstrucción a la dinámica TPK enriquecida. Demuestra
qué memoria necesita: una dodecafase genuina, con posible
\(A_6\ne0\), composición orientada y predecesores torsionales conservados. La
ablación evita atribuir a TPK una limitación que pertenece a una reducción
periódica concreta.

\section{Factorización por el cociente de registro y realización física}

En el alcance de esta sección se declara la relación
\[
\gamma\sim_{\mathcal L}\gamma'
\quad\Longleftrightarrow\quad
\text{\(\gamma,\gamma'\) presentan la misma extensión y }
\mathcal L(\gamma)=\mathcal L(\gamma').
\]
Denotamos por
\(\mathcal E_{108}^{\rm surv}:=
\Gamma_{108}^{\rm enr}/\!\sim_{\mathcal L}\)
el cociente relativo a esta carta. Sea
\[
 \pi_{\rm ext}:\Gamma_{108}^{\rm enr}\longrightarrow
 \mathcal E_{108}^{\rm surv}
\]
la proyección canónica. Por construcción, el extractor es constante en sus
fibras. La afirmación más fuerte según la cual toda carta admisible de una
misma extensión produce automáticamente el mismo registro requeriría demostrar
la compatibilidad entre cartas y no se presupone aquí.

\begin{theorem}[Factorización relativa por el cociente de registro]
\label{thm:descenso-caracter-extension}
Existe una única aplicación
\[
 \overline L:\mathcal E_{108}^{\rm surv}\longrightarrow\ZZ^5
\]
tal que
\[
 \boxed{L_{\rm tick}=\overline L\circ\pi_{\rm ext}.}
\]
Por tanto, \(\chi_{\rm form}\circ\overline L\) es una función de la clase en
el cociente relativo y no del representante de ruta escogido para leerla.
\end{theorem}

\begin{proof}
Si \(\pi_{\rm ext}(\gamma)=\pi_{\rm ext}(\gamma')\), ambas presentaciones
poseen el mismo registro por la relación declarada y, por las fórmulas
anteriores, la misma cinco-tupla. Definimos
\(\overline L([\gamma])=L_{\rm tick}(\gamma)\). La constancia prueba que la
definición es independiente del representante; la sobreyectividad del mapa
al cociente prueba la unicidad.
\end{proof}

La construcción interna llega así hasta
\[
 \Omega_{\HMT}^{\rm surv}
 \longrightarrow\mathcal E_{108}^{\rm surv}
 \xrightarrow{\ \overline L\ }\ZZ^5
 \xrightarrow{\ \chi_{\rm form}\ }
 \operatorname{Fun}(\mathcal P,\RR_{>0}).
\]
El contenido del teorema es positivo y completo en este dominio: una
extensión superviviente determina su clase de registro, la clase determina
una única firma y la firma determina el carácter formal. La especialización
global y la jerarquía completa de torres se construyen en los capítulos
siguientes; las
comparaciones con nombres y unidades físicas se presentan después, sin
alterar esta cadena.
